\section{Protocolo de coordinación TCP y tipos de mensaje}\label{sec:protocol}
El main y los trabajadores se coordinan sobre TCP  con un protocolo orientado a líneas. Cada mensaje es un objeto JSON terminado por salto de línea (NDJSON), de modo que el tráfico resulta legible al depurar. El reparto es por demanda (\emph{work-pulling}), el main mantiene una cola de bloques y entrega uno cada vez que un trabajador lo solicita (un nodo más rápido pide más seguido), y el rendimiento global no queda limitado por el nodo más lento. La \autoref{tab:msgs} lista los tipos de mensaje.

\begin{table}[H]
\centering
\footnotesize
\begin{tabularx}{\linewidth}{l l X}
\hline
\textbf{Mensaje} & \textbf{Dirección} & \textbf{Propósito / campos}\\
\hline
\texttt{HELLO}     & W$\rightarrow$M & registro; id de worker, host, pid, núcleos\\
\texttt{PARAMS}    & M$\rightarrow$W & datos del objetivo, digest, alfabeto, longitud, construcción\\
\texttt{REQUEST}   & W$\rightarrow$M & solicita un bloque del espacio de claves\\
\texttt{ASSIGN}    & M$\rightarrow$W & \{id de bloque, inicio, fin\}\\
\texttt{PROGRESS}  & W$\rightarrow$M & \{bloque, probadas, tasa\}; también sirve de heartbeat\\
\texttt{HEARTBEAT} & W$\rightarrow$M & señal de vida mientras está sin trabajo entre bloques\\
\texttt{DONE}      & W$\rightarrow$M & bloque recorrido por completo sin coincidencia\\
\texttt{FOUND}     & W$\rightarrow$M & \{bloque, clave, índice\} al hallar una coincidencia\\
\texttt{NOWORK}    & M$\rightarrow$W & cola momentáneamente vacía; reintentar\\
\texttt{STOP}      & M$\rightarrow$W & alto global (clave hallada o espacio agotado)\\
\hline
\end{tabularx}
\caption{Protocolo de mensajes maestro--trabajador.}
\label{tab:msgs}
\end{table}

\paragraph{Ciclo de vida.} Un trabajador envía \texttt{HELLO}, recibe \texttt{PARAMS} y alterna \texttt{REQUEST} y \texttt{ASSIGN}, procesando cada bloque en sublotes mientras emite \texttt{PROGRESS}, y reporta \texttt{DONE} o \texttt{FOUND}. Ante un \texttt{FOUND}, el main vuelve a calcular el digest de forma independiente antes de confiar en la clave y solo entonces manda \texttt{STOP} para que ningún nodo siga hasheando. El maestro lanza subprocesos trabajadores locales que se conectan por loopback y crackean junto a los remotos. La reverificación más el \texttt{STOP} garantizan un único resultado confiable.
